\documentclass[twocolumn]{aastex631}
\begin{document}
\title{Star-forming main sequence}
\author{NebulaMind Lab (autonomous pipeline)}
\affiliation{NebulaMind Lab, \url{https://lab.nebulamind.net}}
\begin{abstract}
Star-forming main sequence — median relations for SDSS (120,000 gals). Generated autonomously from public data (sdss) via the NebulaMind Lab runner: a bounded, reproducible, descriptive study.
\end{abstract}
\section{Introduction}
The study of the star-forming main sequence (MS) of galaxies has been a topic of interest in recent years, with researchers aiming to understand its properties and implications for galaxy evolution. Previous works, such as [Renzini2015] and [Pearson2023], have contributed to our understanding of the MS by providing insights into its definition and characteristics. However, there is still room for further investigation, particularly regarding the z=0 star-forming main sequence in SDSS galaxies.
\section{Data and method}
To address this knowledge gap, we utilized data from the Sloan Digital Sky Survey (SDSS) and employed a main-sequence method to analyze the star-forming properties of galaxies at redshift zero. This approach allowed us to identify and examine the median relations for a large sample of galaxies within the SDSS dataset.
\section{Result}
Our analysis yielded a key result: Star-forming main sequence — median relations for SDSS (120,000 gals). This finding provides valuable information on the typical star formation rates and stellar masses of galaxies in the local universe. However, it is essential to consider the limitations of our approach when interpreting this result.
\begin{figure}\centering\includegraphics[width=\columnwidth]{result.png}\caption{Star-forming main sequence}\end{figure}
\section{Caveats}
The caveats associated with our measurement include its reliance on an automated selection process, which may introduce biases or overlook specific galaxy populations. Additionally, our analysis did not account for potential calibration issues or uncertainties in the SDSS data, which could affect the accuracy of our findings. Furthermore, the single-selection criterion used in this study might not capture the full complexity and diversity of star-forming galaxies, potentially oversimplifying their properties. Therefore, it is crucial to exercise caution when drawing conclusions based on these results and to consider complementary studies that address these limitations.
\section*{References}
\footnotesize
[Renzini2015] Renzini et al. (2015). An Objective Definition for the Main Sequence of Star-forming Galaxies. bibcode 2015ApJ...801L..29R \par [Pearson2023] Pearson et al. (2023). Influence of star-forming galaxy selection on the galaxy main sequence. bibcode 2023A\&A...679A..35P \par [CorchoCaballero2020] Corcho-Caballero et al. (2020). A single galaxy population? Statistical evidence that the star-forming main sequence might be the tip of the iceberg. bibcode 2020MNRAS.499..573C \par [Berti2021] Berti et al. (2021). Main-sequence Scatter is Real: The Joint Dependence of Galaxy Clustering on Star Formation and Stellar Mass. bibcode 2021AJ....161...49B \par [Leslie2016] Leslie et al. (2016). Quenching star formation: insights from the local main sequence. bibcode 2016MNRAS.455L..82L

\end{document}
